%%%%%%%%%%%%%%%%%%%%%%%%%%%%%%%%%%%
% Main Text
%%%%%%%%%%%%%%%%%%%%%%%%%%%%%%%%%%%

\section{Introduction}\label{sec:intro}

Asian Americans represent the fastest-growing racial group in the United States, yet their experiences with discrimination and identity formation remain underexplored in economic research.\footnote{The 2020 Census counted more than 20 million Asian Americans---6.4 percent of the population---nearly double the number counted two decades earlier. The Asian American population numbers are based on the author's calculations from the Current Population Survey and US Census data \autocite{floodsarahIntegratedPublicUse2021a}.} Unlike other minority groups, Asian Americans occupy a distinctive position in America's racial hierarchy---simultaneously experiencing discrimination and being labeled as ``perpetual foreigners'' while also being characterized through the ``model minority'' stereotype \autocite{foukaImmigrantsAmericansRace2022}. This dual status creates incentives around racial identity choices that differ from other groups' experiences. Whether those choices respond to local prejudice is a question with direct consequences for policy, because nearly every policy that counts, protects, or monitors Asian Americans does so through self-reported race. The federal standards for collecting data on race and ethnicity, revised in 2024, govern the statistics used to enforce civil rights law, allocate federal funds, and track disparities \autocite{ombRevisionsStatisticalPolicy2024}, and the COVID-19 Hate Crimes Act of 2021 directed federal and state agencies to improve the reporting of anti-Asian hate crimes \autocite{covid19HateCrimesAct2021}. If people of Asian ancestry are less likely to identify as Asian where prejudice is stronger, then the population counts, outcome gaps, and victimization rates that these policies rely on are measured with an error that is itself a function of prejudice.

An extensive literature has documented how Asians and Whites compare across multiple outcomes, yet the role of identity selection in shaping the Asian-White gaps remains understudied \autocite{chiswick1983analysis, duleep2012economic, hilger2016upward, arabsheibani2010asian}. The challenge lies in defining and measuring racial identity, particularly when individuals possess agency in how they racially self-identify. If reporting Asian racial identity represents a strategic choice influenced by local discrimination, measured gaps may systematically vary across geographic contexts in ways that previous research has not fully explored.

Various contextual factors, including anti-Asian sentiment and stereotypes, can influence how individuals navigate their racial identity choices. The COVID-19 pandemic has brought renewed attention to how external hostility shapes Asian American experiences \autocite{gover2020anti}. In this paper, I examine the determinants of Asian racial identity reporting and analyze how Asian Americans strategically select between Asian and White racial identities. Specifically, I investigate how anti-Asian bias, education, and family income shape decisions to identify racially as Asian American. I also break down the analysis by generation and family structure (interracial versus endogamous parents) and examine what racial identities individuals with objective Asian ancestry report (Asian only, White only, multiracial, etc.).

I explore how individual characteristics and societal attitudes toward Asian Americans influence racial identity reporting. I use identity and ancestry data from the Current Population Survey (CPS) combined with measures of anti-Asian bias derived from Harvard's Project Implicit Association Test (IAT), the American National Election Studies (ANES), and hate crimes targeting Asian Americans.\footnote{The IAT data comes from Harvard's Project Implicit \autocite{greenwaldMeasuringIndividualDifferences1998}. Implicit bias measures have gained prominence in economics, with IAT scores correlating with economic outcomes \autocite{chettyRaceEconomicOpportunity2020,gloverDiscriminationSelfFulfillingProphecy2017}, voting patterns \autocite{friesePredictingVotingBehavior2007}, and health disparities \autocite{leitnerRacialBiasAssociated2016}.} I ground my analysis in the theoretical framework of \textcite{akerlofEconomicsIdentity2000}, explicitly modeling how external prejudice creates differential utility from identity choices and establishing conditions under which individuals strategically modify their racial self-presentation. I also use CPS data to study the effect of education (parental and individual) and family income on Asian racial identity reporting.

Measuring identity choices outside of lab settings is challenging, requiring both objective ancestry indicators and subjective identity measures. I leverage birthplace and ancestry data to construct objective Asian ancestry measures, then analyze deviations between objective ancestry and subjective racial identity. I find that racial identity reporting responds to both individual characteristics and environmental factors reflecting local discrimination levels. 

I find substantial Asian racial identity attrition across generations, driven primarily by interracial family structures. While children of endogamous Asian parents report Asian identity at rates above 90\%, only about a third of second-generation children with one Asian and one White parent do the same. Using multinomial logit models, I show that anti-Asian bias shifts identity choices away from ``Asian only'' and toward ``White only,'' with effects that are larger in families where the father is Asian and that scale with the number of Asian grandparents among the third generation. On average, parental education increases the probability of Asian identification. However, this relationship reverses in high-bias states, where individuals with higher parental education are less likely to identify as Asian than those with lower parental education under the same levels of bias. Because of this shifting composition, the Asian-identified population has a lower average socioeconomic status in high-bias areas. Consequently, outcomes measured on self-reported race understate the outcomes of people of Asian ancestry where prejudice is strongest, and conventional estimates of Asian--White disparities are biased there.

The policy stakes follow from the direction of this selection. First, the Asian population that official statistics count falls short of the population of Asian ancestry by the most where anti-Asian prejudice is strongest, so the number of people exposed to that prejudice is understated where protection matters most. Second, because the children who stop identifying as Asian in high-bias states are disproportionately those with college-educated mothers, the Asian-identified population that remains there has lower socioeconomic status than the population of Asian ancestry. Outcomes measured on self-reported race therefore understate the outcomes of people of Asian ancestry in high-bias states, and the comparisons of Asian and White outcomes that guide civil rights enforcement, the targeting of programs by race, and the assessment of immigrant integration are distorted in a way that varies systematically with prejudice. \textcite{duncanIdentifyingLaterGenerationDescendants2018} show that selective ethnic attrition of this kind biases measures of intergenerational progress for Hispanics, and the National Academies raised the same concern for integration statistics generally \autocite{watersIntegrationImmigrantsAmerican2015}; the estimates in this paper put magnitudes on the mechanism for Asian Americans and show which families drive it.

This research contributes to several bodies of literature. First, it extends the economics of identity framework \autocite{akerlofEconomicsIdentity2000} and stratification economics \autocite{darityEconomicsIdentityOrigin2006,darityPositionPossessionsStratification2022} by examining how racial stereotypes---both positive and negative---influence identity choices among Asian Americans, who face a unique utility landscape where identity can simultaneously signal competence and foreignness \autocite{charnessSocialIdentityGroup2020,atkinHowWeChoose2021,goldsmithDarkLightSkin2007,hamiltonSheddingLightMarriage2009,dietteSkinShadeStratification2015}. Qualitative research documents how individuals strategically adapt racial presentations across contexts \autocite{waters1990ethnic,telles2008generations,zhou1997segmented,cornell2006ethnicity}; this paper quantifies those choices by comparing objective ancestry with subjective identification across varying bias environments.

Second, it extends research on discrimination \autocite{bertrandAreEmilyGreg2004,charlesPrejudiceWagesEmpirical2008,bursztynImmigrantNextDoor2022,bordaloStereotypes2016} by examining how prejudice influences racial self-identification itself. Within immigration research, the paper builds on studies of assimilation \autocite{abramitzkyCulturalAssimilationAge2016,abramitzkyNationImmigrantsAssimilation2014}, but unlike European immigrant groups, Asian Americans face persistent ``perpetual foreigner'' stereotypes that complicate integration regardless of generational status \autocite{foukaImmigrantsAmericansRace2022}, and the model minority myth means Asian racial identity may carry both benefits and costs depending on context \autocite{mengIntermarriageEconomicAssimilation2005}.

This paper most closely relates to recent economic research on racial identity fluidity \autocite{hadah2024hispanicidentity, antmanEthnicAttritionObserved2016,antmanIncentivesIdentifyRacial2015,antmanAmericanIndianCasinos2021}. While previous work focused primarily on Hispanic ethnic attrition, Asian American identity choices operate through different mechanisms due to distinct stereotypes, discrimination patterns, and socioeconomic profiles.

\section{Theoretical Framework}\label{sec:model}

I develop a theoretical framework for understanding racial identity choice that extends \textcite{akerlofEconomicsIdentity2000} to incorporate stereotype-specific costs and benefits. Unlike generic minority identification models, this framework recognizes that Asian Americans face unique utility trade-offs where racial identity can signal both positive attributes (academic achievement, work ethic) and negative characteristics (foreignness, social exclusion).

Formally, individual $i$ belongs to racial group $r_i \in \{A, W\}$, where $A$ represents Asian and $W$ represents White. Agent $i$'s utility depends on their actions and how those actions interact with their chosen racial identity $I_i$:

\begin{equation}
U_i = U_i(\pmb{a_i}, \pmb{a_{-i}}, I_i)\label{eq:util}
\end{equation}

Individual identity $I_i$ reflects personal actions, others' behaviors toward them, and societal expectations associated with their racial group:

\begin{equation}
I_i = I_i(\pmb{a_i}, \pmb{a_{-i}}; \pmb{S}_{r_{i}})\label{eq:identity}
\end{equation}

where $\pmb{a_i}$ represents individual $i$'s actions, $\pmb{a_{-i}}$ captures others' actions affecting $i$'s identity (including anti-Asian bias), and $\pmb{S}_{r_{i}}$ denotes societal stereotypes and expectations associated with racial group membership.\footnote{This extends \textcite{akerlofEconomicsIdentity2000}'s prescription concept to encompass both negative stereotypes and positive model minority expectations.}

The key insight is that $\pmb{S}_{A}$ includes both positive stereotypes (academic excellence) and negative ones (perpetual foreigner status). Thus creating context-dependent utility from Asian racial identification.

Individual $i$ selects actions $\pmb{a_i}$ to maximize utility given their racial group $r_i$, associated stereotypes $\pmb{S}_{r_{i}}$, and others' actions $\pmb{a_{-i}}$. The first-order condition becomes:

\begin{equation}
\frac{\partial U_i}{\partial \pmb{a_i}} + \frac{\partial U_i}{\partial I_i} \cdot \frac{d I_i}{d \pmb{a_i}} = 0\label{eq:foc}
\end{equation}

The solution $a_{i}^{\star}$ yields utility $U_{i}^{\star}$. Now suppose individuals can strategically choose their racial identity at cost $c$. They will switch identities when $\tilde{U_{i}}^{\star} \geq U_{i}^{\star} + c$, where $\tilde{U_{i}}^{\star}$ represents utility under the alternative racial identity.

Identity switching occurs when benefits $\tilde{U_{i}}^{\star} - U_{i}^{\star}$ exceed costs $c$. These net benefits are non-zero only when $\frac{d I_i}{d \pmb{a_i}} \neq 0$ and $\frac{\partial U_i}{\partial I_i} \neq 0$. This framework suggests empirical analysis should focus on: (1) individual characteristics affecting optimal actions under different racial identities, (2) contextual factors (anti-Asian bias) creating differential treatment by racial group, (3) populations with low switching costs $c$, and (4) groups whose utility significantly depends on racial identity.

The model predicts that anti-Asian bias increases the utility differential between White and Asian racial identity, making identity switching more attractive. Mixed-race individuals face lower switching costs due to phenotypic ambiguity, while later-generation Asian Americans may find identity switching more feasible due to cultural assimilation.

\section{Data Sources and Measurement Strategy}\label{sec:data}

In this section, I describe the datasets I use in the analysis. To examine relationships between social attitudes and Asian racial identity reporting, I need both subjective and objective Asian identity measures for identifying appropriate Asian American subgroups. I use the Integrated Public Use Microdata Series (IPUMS) Current Population Survey (CPS) \autocite{floodsarahIntegratedPublicUse2021a} with ancestry information through the places of birth of individuals, their parents, and grandparents to construct objective identity measures. I develop composite anti-Asian bias measures using \textcite{lubotskyInterpretationRegressionsMultiple2006}'s multiple proxy regression method to reduce attenuation bias.

\subsection{Measuring Asian Racial Identity}\label{subsec:cps}

I measure Asian racial identity using CPS data from 2004--2021. Following \textcite{antmanEthnicAttritionObserved2016,antmanEthnicAttritionAssimilation2020}, I utilize birthplace information for individuals, parents, and grandparents to create objective Asian ancestry indicators for minors under age 17 living with parents.\footnote{This approach parallels previous research but focuses on racial rather than ethnic categorization.} The methodology allows for the identification of first-, second-, and third-generation Asian Americans (see Figure \ref{fig:diag} for visual representation). I also use a CPS sample of adults aged 18+ to examine whether their racial identity reporting patterns differ from minors; for the adult sample, I can only identify first- and second-generation.\footnote{Since adults still living with parents might not be representative of the overall adult population and is a rare occurrence, I do not link adults to their parents.}

The objective ancestry measure---distinct from subjective racial identification where respondents select ``Asian'' as their race---depends on birthplaces across three generations.\footnote{For this analysis, Asian countries comprise East Asian and Southeast Asian nations, including China, Hong Kong, Taiwan, Japan, Korea, Mongolia, Cambodia, Indonesia, Laos, Malaysia, Philippines, Singapore, Thailand, and Vietnam, but exclude South Asian and Middle Eastern countries, consistent with standard demographic classifications.} The three identifiable generations include: 1) first-generation immigrants born in Asian countries with both parents also born in Asian countries, 2) second-generation individuals who are US-born citizens with at least one parent born in an Asian country, and 3) third-generation Asian Americans who are US-born citizens with two US-born parents and at least one grandparent born in an Asian country. 

While the ancestry measure provides an objective assessment of Asian heritage, it may not capture all nuances of racial identity. For instance, White individuals with Asian ancestry born to non-Asian parents in Asia, such as on American military bases, may be classified as Asian in the data. To avoid potential misclassification, I remove individuals who report that they were born abroad of American parents. The final sample includes Asian Americans, first-, second-, and third-generation immigrants aged 17 and younger living with parents between 2004 and 2021. I present the summary statistics in Table (\ref{tab:sumstat1}). The adult sample includes first- and second-generation Asian American immigrants aged 18 and older between 2004 and 2021. I show the summary statistics for the adult sample in Table (\ref{tab:sumstat-adults}).

CPS relies on household respondents (parents or caregivers) to report children's racial identity. Asian identity reporting rates are consistent regardless of whether mother (72\%), father (72\%), or child/other caregiver (87\%) serves as respondent (Table \ref{tab:hispbyproxy}), and attrition patterns among adults align with those observed in children (Tables \ref{tab:asianbygen} and \ref{tab:asianbygen-adults}). I address proxy reporting concerns more fully in Section~\ref{sec:robcheck}.

The overall sample comprises 49\% females, with 65\% self-reporting Asian racial identity---answering affirmatively to ``what is your race.'' Average age is 8.4 years. Approximately 52\% of mothers and 52\% of fathers hold college degrees. Additional summary statistics for the overall sample and each generation appear in Table (\ref{tab:sumstat1}).

Using parental and grandparental birthplaces, I objectively identify ethnic ancestry and categorize different family types. For second-generation children, parental birthplaces create three objective categories: Objectively Asian-father-Asian-mother (AA), Objectively Asian-father-White-mother (AW), and Objectively White-father-Asian-mother (WA)

Similarly, grandparental birthplaces create 15 objective categories for third-generation children: (1) objectively Asian paternal grandfather-Asian paternal grandmother-Asian maternal grandfather-Asian maternal grandmother (AAAA); (2) objectively White paternal grandfather-Asian paternal grandmother-Asian maternal grandfather-Asian maternal grandmother (WAAA); (3) objectively Asian paternal grandfather-White paternal grandmother-Asian maternal grandfather-Asian maternal grandmother (AWAA), etc.

My analysis employs a subsample of the US population; Tables \ref{tab:asianbygen} and \ref{tab:sumstat-adults} demonstrate sufficient observations across generations for both adults and children. Consistent with literature on ethnic and racial identity fluidity among Asian and Hispanic Americans, I document significant attrition among third-generation Asian Americans.\footnote{\textcite{duncanIdentifyingLaterGenerationDescendants2018,duncanSocioeconomicIntegrationImmigrant2018, antmanEthnicAttritionObserved2016,antmanEthnicAttritionAssimilation2020} document substantial identity attrition among various groups.}

\subsection{Measuring Anti-Asian Sentiment}\label{sub:lw-bias}

I construct anti-Asian sentiment measures using implicit association tests, American National Election Studies, and hate crimes targeting Asian Americans from 2004--2021. 

The implicit association test measures conceptual associations---for example, linking Asian Americans with negative stereotypes---and evaluative responses. Respondents rapidly categorize words into screen-displayed categories. Figure (\ref{fig:iatexamples}) shows examples from Harvard's Project Implicit skin tone test.

I employ Asian-focused IAT data to construct anti-Asian prejudice proxies \autocite{greenwaldMeasuringIndividualDifferences1998}. The IAT measures bias direction and magnitude while capturing unconscious biases individuals may be unwilling to report, and previous research demonstrates the difficulty of manipulating IAT scores \autocite{egloffPredictiveValidityImplicit2002}. A meta-analysis by \textcite{greenwaldMeasuringIndividualDifferences1998} finds significantly higher predictive validity for IAT compared to self-report measures.\footnote{IAT participation is voluntary, potentially creating selection bias. However, IAT-reflected bias serves as a proxy for prejudiced attitudes \autocite{chettyRaceEconomicOpportunity2020}.} 

However, some research questions IAT predictive validity claims. Implicit Association Tests may not reliably measure or predict implicit prejudice or biased behaviors. Research shows implicit biases experience minor, temporary intervention-induced changes. Additionally, implicit bias fails to predict dictator game contributions or social pressure susceptibility, highlighting distinctions between implicit bias and biased actions \autocite{arkesAttributionsImplicitPrejudice2004,forscherMetaanalysisProceduresChange2019,leeDoesImplicitBias2018}. Therefore, I supplement IAT with explicit bias measures from American National Election Studies (ANES) and hate crimes against Asian individuals to construct a composite bias measure.

I develop a second proxy from the American National Election Studies (ANES), averaging responses across several 2004--2020 questions measuring racial animus \autocite{anes2021}.\footnote{Questions parallel those used by \textcite{charlesPrejudiceWagesEmpirical2008}: (1) ``Conditions Make it Difficult for Blacks to Succeed'', (2) ``Blacks Should Not Have Special Favors to Succeed'', (3) ``Blacks Must Try Harder to Succeed'', (4) ``Blacks Gotten Less than They Deserve Over the Past Few Years'', and (5) ``Feeling Thermometer Toward Asians.''} While the ANES racial animus questions primarily focus on attitudes toward Black Americans, research demonstrates that racial prejudices are highly correlated across different minority groups, with individuals who express bias toward one racial minority typically holding similar attitudes toward others \autocite{almasalkhi2023links, mora2020antiblackness}. These measures therefore capture broader patterns of racial animus that extend beyond anti-Black sentiment specifically. When combined with the Asian-focused IAT measures and hate crimes against Asian Americans in my composite index, this multi-proxy approach weights the bias measure more heavily toward Asian-specific prejudice while still capturing the general racial climate. Finally, I incorporate Uniform Crime Reports (UCR) data quantifying hate crimes against Asian Americans \autocite{ucrbook}.

To reduce attenuation bias and measurement error, I follow \textcite{lubotskyInterpretationRegressionsMultiple2006} in constructing composite bias measures using IAT, ANES racial animus measures, and hate crimes against Asian Americans.\footnote{Additional methodological details appear in the Online Appendix.} Figure (\ref{fig:skiniat}) graphically represents bias measures over time in the most and least biased locations. Figure (\ref{fig:Asian-twostates}) shows Asian racial identity reporting in the two most and least biased locations. Lower scores indicate less bias; higher scores indicate greater racial animus. The composite is scaled in the units of the IAT $D$ score, a standardized difference in response latencies, and throughout the paper I refer to a one-unit change in the index as a one standard deviation change in this sense. Relative to the variation in the data, one unit is a large change. The index is measured in two-year blocks (2004--05, 2008--09, 2012--13, 2016--17, and 2020--21), so each state contributes at most five distinct values, and the panel is unbalanced: 29 states enter in 2004, 11 in 2008, and 11 in 2012. Across the 432 state-year cells with a bias measure, the index has a standard deviation of 0.42, of which the between-state component (the standard deviation of state means) is 0.18 and the within-state component over time is 0.39; within a given year, the standard deviation across states averages 0.26. A one-unit increase, the change used in the tables and figures, is therefore about 2.4 standard deviations of the state-year distribution, or the gap between the least biased places (Washington, DC, and Vermont) and the most biased (North Dakota) in 2020. Appendix Figures~\ref{fig:skiniat-maps} and~\ref{fig:iat-map-all} map the index by year and averaged over 2004--2021.

\section{From the Data: Asian Racial Identity and Attrition}\label{sec:attrition}

Table~\ref{tab:asianbygen} reports the distribution of self-reported racial identity among children of Asian ancestry by generation and ancestry, using the three categories that the multinomial logit models below distinguish (``Asian only,'' ``White only,'' and ``Asian and White'') plus a residual category.\footnote{Beginning in 2003, the CPS allowed respondents to report multiple racial identities. ``Asian and White'' combines any response that includes both Asian and White, with or without Pacific Islander; ``Other'' collects Asian or Pacific Islander, other single races, and multiracial combinations that do not include both Asian and White. Appendix Figures~\ref{fig:histogram-all}--\ref{fig:histogram-thirdgen} show a finer six-way breakdown.} Among first-generation children, 94\% report ``Asian only.'' The share falls to 72\% in the second generation and 29\% in the third. Attrition is concentrated in interracial families. Among second-generation children with two Asian-born parents, 96\% report ``Asian only,'' compared with 31\% of children with one Asian-born parent and one parent born outside Asia (86\% of whom are US-born); the rest of the latter group split between ``White only'' (25\%) and ``Asian and White'' (37\%). The split differs by which parent is Asian: 36\% of children of an Asian father and a White mother report ``Asian only'' versus 29\% of children of a White father and an Asian mother, who more often report ``Asian and White'' (40\% versus 29\%).

Among third-generation children, the number of Asian-born grandparents predicts identity almost completely. Children with four Asian grandparents report ``Asian only'' at 93\%; with three, 74\%; with two, 24\%; and with one, 8\%, at which point ``White only'' is the modal response (54\%). Adults show the same pattern (Table~\ref{tab:asianbygen-adults}): 93\% of second-generation adults with two Asian-born parents report ``Asian only'' against 33\% of those with one Asian-born parent, and among the latter ``White only'' (34\%) is as common as ``Asian only.''

These patterns show that racial identity becomes more fluid across generations. Endogamous Asian families maintain high rates of Asian identification, while interracial family structures leave most children choosing between a multiracial and a White-only identity, which is the margin on which bias and family background can operate.

\section{Empirical Approach}\label{sec:empstrat}

To understand associations between Asian racial self-identification and anti-Asian bias, I estimate a multinomial logit model in which the dependent variable $Y_{ist}^g$ captures categorical racial identity choices: (1) Asian only, (2) White only, (3) Asian and White/Pacific Islander:

\begin{align}
\log\left(\frac{P(Y_{ist}^g = j)}{P(Y_{ist}^g = \text{Asian only})}\right) &= \beta_{1j}^g AntiAsianBias_{st} + \beta_{2j}^g DadCollegeGrad_{ist}\nonumber \\ 
& + \beta_{3j}^g MomCollegeGrad_{ist} + \beta_{4j}^g Female_{ist} + X_{ist}^g\pi_j  \nonumber \\ 
& + \gamma_{rtj} + \varepsilon_{istj}; \quad j \in \{\text{White only}, \text{Asian and White}\} \label{eq:multinomial_logit}
\end{align}

where $Y_{ist}$ represents the categorical racial identity of person $i$ in state $s$ at interview time $t$, $j$ indexes the identity categories, with ``Asian only'' as the reference category, $P(Y_{ist}^g = j)$ denotes the probability of person $i$ in state $s$ at time $t$ choosing identity $j$, $\mathbf{X}_{ist}^g$ represents the vector of controls that include quartic age, Asian population fraction, parent types, grandparent types, and generation dummies, $\boldsymbol{\beta}_j^g$ denotes the coefficient vector for outcome $j$ and generation $g$, and $J$ indexes the three identity categories. Finally, $\gamma_{rtj}$ represents region-time fixed effects that would control for region $\times$ year specific shocks affecting identity choice $j$, and $\varepsilon_{istj}$ is the error term.

The model specification allows the effects of anti-Asian bias and other covariates to vary across identity choices. For instance, $\beta_{1,\text{White only}}^g$ captures how anti-Asian bias affects the log-odds of choosing ``White only'' versus ``Asian only'' identity, while $\beta_{1,\text{Asian and White}}^g$ measures the bias effect on choosing ``Asian and White'' versus ``Asian only.'' This framework enables analysis of how bias, sex, and parental education differentially influence various identity strategies available to individuals with Asian ancestry.

The coefficients of interest are $\beta_{1j}^g$, $\beta_{2j}^g$, $\beta_{3j}^g$, and $\beta_{4j}^g$, which capture how anti-Asian bias, parental education, and sex influence the likelihood of selecting each identity category relative to ``Asian only.'' I estimate separate models for each generation $g \in \{1,2,3\}$ to assess whether these relationships differ across generational status.

The multinomial specification is appropriate because choosing ``White only'' identity represents a fundamentally different strategy than selecting ``Asian and White'' multiracial identity, a distinction particularly important for mixed-ancestry individuals who may strategically choose between complete ethnic distancing and maintaining partial ethnic connection.

Because multinomial logit coefficients are log-odds ratios relative to the ``Asian only'' category, I report average marginal effects on the probability of each identity. For anti-Asian bias, the marginal effect for each child is the derivative of the predicted probability of an outcome with respect to the bias index, evaluated at that child's own covariate values; I average these derivatives, without weights, over the estimation sample, so the reported number is the average change in probability per one-unit (one standard deviation) increase in bias.\footnote{The derivative is computed numerically as a central difference with a step of one-tenth of a standard deviation of the index; the perturbation is truncated at the sample range of the index, which pulls the derivatives for children in the most and least biased cells slightly toward zero.} For the binary covariates (female, mother college graduate, father college graduate), the marginal effect is the average over the sample of the change in predicted probability from switching the indicator from 0 to 1. For inference, I draw 1{,}000 bootstrap samples of children with replacement, stratified by region $\times$ year cell, re-estimate the model in each, and recompute the average marginal effects. The point estimates I report are the medians of the 1{,}000 bootstrap estimates, and the confidence intervals are their 2.5th and 97.5th percentiles. Because the bootstrap resamples children rather than states, these intervals reflect sampling variation in the individual data and do not account for dependence among children in the same state-year cell. That dependence matters: in the pooled model of Appendix Table~\ref{tab:pooled-mnl-bias-x-aw}, standard errors on the bias effects clustered at the state level are about three times the conventional standard errors that treat children as independent, so the bootstrap intervals in the figures should be read as understating uncertainty by a similar factor. I therefore report the pooled bias effect and the key AW--WA contrast from that model, with state-clustered standard errors, alongside the bootstrap estimates.

\section{Results and Robustness Checks}\label{sec:results}

Figure~\ref{fig:marginal-effects-pooled} presents average marginal effects from the multinomial logit models for the pooled mixed-race second-generation sample (panel a) and the pooled third-generation sample (panel b), the groups for which identity is fluid. Appendix Figures~\ref{fig:marginal-effects-second-parental}, \ref{fig:marginal-effects-second-adults-pooled}, and~\ref{fig:marginal-effects-third-grandparental} report the same estimates separately by parental configuration, for adults, and by number of Asian grandparents.

\textbf{Second-Generation Subsample.} Among children of mixed-race parents (Figure~\ref{fig:marginal-effects-second-pooled}), a one standard deviation increase in anti-Asian bias increases ``White only'' identity by 12 percentage points (pp) and decreases ``Asian only'' by 10~pp. A pooled model with a common bias effect and standard errors clustered at the state level (Appendix Table~\ref{tab:pooled-mnl-bias-x-aw}, Panel B) gives a similar magnitude, $+$14.6~pp (standard error 5.7) for ``White only,'' with the offsetting decline falling almost entirely on ``Asian only'' (12.3~pp) rather than ``Asian and White'' (2.4~pp), so the result does not depend on the individual-level bootstrap. Maternal college education increases ``Asian only'' ($+$5~pp) and decreases ``White only'' ($-$6~pp). Paternal college education increases ``Asian and White'' ($+$5~pp) and decreases ``White only'' ($-$4~pp). Results among second-generation adults are directionally similar (Appendix Figure~\ref{fig:marginal-effects-second-adults-pooled}).

I also estimate the model separately by parental composition---Asian father--White mother (AW) versus White father--Asian mother (WA) families (Appendix Figure~\ref{fig:marginal-effects-second-parental}). Point estimates suggest that bias effects on ``White only'' identity are larger in AW families ($+$20~pp) than in WA families ($+$10~pp), and that maternal college education effects differ in direction and magnitude across the two family types.

Comparing these two sets of estimates requires care, because they come from separately estimated models. A difference in average marginal effects reflects differences in coefficients but also differences in the two groups' covariate distributions, and the coefficients themselves are not comparable across multinomial logits whose error scales are normalized separately. I therefore test the AW--WA difference in two ways. Appendix Table~\ref{tab:between-model-second} and Appendix Figure~\ref{fig:between-model-second} report between-model tests of the differences in average marginal effects, with bootstrap standard errors for each difference. Appendix Table~\ref{tab:pooled-mnl-bias-x-aw} instead pools the two groups in a single multinomial logit that constrains every coefficient to be equal across AW and WA families except an interaction between bias and an AW indicator, and computes the average marginal effect of bias for both groups over the same pooled sample, so that the comparison holds the covariate distribution fixed. Both approaches agree on the main difference. In the pooled model, the bias $\times$ AW interaction in the ``White only'' equation is 0.89 (standard error 0.32), and a Wald test with state-clustered standard errors rejects equal bias effects in the two equations ($\chi^2(2) = 10.5$, $p = 0.005$); a one standard deviation increase in bias raises ``White only'' identification by 24.9~pp in AW families and 9.9~pp in WA families, a difference of 15.1~pp (standard error 4.2). The corresponding difference for ``Asian only'' is $-$10.5~pp (standard error 6.4, $p = 0.10$), and for ``Asian and White'' it is not distinguishable from zero.\footnote{A linear probability model for the binary Asian indicator with the same pooled specification gives a bias $\times$ AW interaction of $-$0.08 (standard error 0.06).} The effect of bias is therefore stronger when the father is Asian, pushing children toward a ``White only'' identity. The between-model tests also find that maternal and paternal college education effects differ across family types for ``Asian only'' and ``White only''; the remaining AW--WA differences are not statistically distinguishable. The between-model tests rest on the individual-level bootstrap and inherit the caveat from Section~\ref{sec:empstrat}: their standard errors do not account for state-level dependence and are probably about a third of what state clustering would give, so differences that are significant only by that criterion should be read as suggestive. The AW--WA difference in the bias effect is the one that survives state-clustered inference.

\textbf{Third-Generation Subsample.} Among pooled third-generation children (Figure~\ref{fig:marginal-effects-third-pooled}), bias effects on identity are not statistically significant; the one exception is paternal college education, which increases ``Asian and White'' identification by 8~pp. Because bias and parental education should matter most for those with the strongest ancestral ties, I split the sample by number of Asian grandparents (Appendix Figure~\ref{fig:marginal-effects-third-grandparental}). Both sets of effects grow with Asian ancestry: the effect of bias on ``White only'' identity is negligible among children with one Asian grandparent but reaches $+$19~pp among those with three, and the effect of maternal college education on ``Asian only'' grows to $+$32~pp over the same gradient. The between-model tests I conduct (Appendix Figure~\ref{fig:between-model-third} and Appendix Table~\ref{tab:between-model-third}) find these differences statistically significant by the bootstrap criterion, subject to the same caveat: comparing one versus three Asian grandparents, the bias effect on ``White only'' and both parents' college effects on ``Asian only'' and ``White only'' differ significantly. Maternal college effects differ between one and two grandparents, and bias and parental education effects differ between two and three. Gender effects never differ across groups.

The effects of bias and parental education on identity are heterogeneous across ancestry types: even where cross-group differences fall short of statistical significance, point estimates increase steadily with the number of Asian grandparents. Consistent with the theoretical framework, individuals with lower switching costs---those with fewer Asian grandparents and greater racial ambiguity---show smaller marginal effects of bias because they can switch identities even at low levels of bias, while those with more Asian ancestry require stronger incentives to overcome higher switching costs.\footnote{\textcite{hadah2024hispanicidentity} finds similar effects on third-generation Hispanic immigrants.}

\subsection*{The interaction between state-level bias and socioeconomic status}

To assess whether bias and socioeconomic status jointly affect Asian identification, I estimate linear probability models that interact demeaned state-level bias with sex and with parental education (own education and family income for adults). Table~\ref{tab:heterogeneity-interactions} reports the interaction coefficients for mixed-race second-generation children (Panel A), mixed-race second-generation adults (Panel B), and third-generation children by number of Asian grandparents (Panel C), separately by family type and pooled.
 
Education amplifies the negative effect of bias on Asian identification, and the pattern is concentrated in families with an Asian father. Among second-generation AW children, the interaction between bias and maternal college education is $-$0.23 (standard error 0.10): for a one standard deviation increase in bias, children of college-educated mothers lose 23 percentage points more Asian identification than children of mothers without a degree. AW adults show the same pattern with their own college education ($-$0.08, standard error 0.05, significant at the ten percent level). In WA families no interaction is distinguishable from zero in either the children or the adult sample. Pooling AW and WA families, the bias-by-maternal-college interaction is $-$0.07 (standard error 0.03) among children and not distinguishable from zero among adults. Among third-generation children, the point estimates are negative and grow with the number of Asian grandparents ($-$0.02, $-$0.08, and $-$0.30 for one, two, and three grandparents), but the subsamples are small and none is precisely estimated; the pooled third-generation interaction is $-$0.07 (standard error 0.03). Interactions with sex, with paternal education, and with family income are not distinguishable from zero in any panel.
 
These results imply that higher-educated individuals of Asian ancestry are disproportionately likely to exit Asian identification in high-bias states. The remaining Asian-identified population is therefore negatively selected on education and socioeconomic status, so outcomes measured on self-reported race understate the outcomes of people of Asian ancestry in states where anti-Asian prejudice is strongest.

\subsection{Robustness Checks}\label{sec:robcheck}

I examine four potential threats to identification: endogenous interracial marriage, selective migration, proxy reporting, and fixed state characteristics. If anti-Asian bias influences interracial marriage rates, it could affect the composition of the mixed-race sample. I regress an indicator for an interracial couple on state-level bias with partner-specific controls among the parents of second-generation children (Table~\ref{regtab-logit-02}). The point estimates are small. A one standard deviation increase in bias is associated with a 4~pp higher probability that a child's parents are an interracial couple (standard error 3~pp), a 3~pp increase among children of Asian-born mothers, and a 1~pp decrease among children of Asian-born fathers, and none of the three is statistically distinguishable from zero with standard errors clustered at the state level.\footnote{The previous version of this table clustered standard errors at the household level, which treats the state-year bias measure as if it varied across households and understated the standard errors by a factor of about three; the point estimates are unchanged.} The opposite signs for Asian-born mothers and fathers deserve comment, because in a closed marriage market every Asian woman who marries a White man is matched by an Asian man who does not marry an Asian woman. This sample is not a closed market. It conditions on couples with a co-resident child, it includes only couples in which one partner was born in Asia and the other was not, and Asian-born women are in such couples far more often than Asian-born men (20,247 children of a White father and an Asian mother against 9,055 children of an Asian father and a White mother). A bias-induced change in whether Asian men marry at all, or in whether they form a family with a co-resident child, would appear as a difference between the two columns rather than as an interracial marriage. The estimates are therefore consistent with bias affecting marriage itself for Asian men, not only partner choice, although their imprecision means the data do not settle this.

Using 2004--2021 American Community Survey data \autocite{floodsarahIntegratedPublicUse2021}, I test whether anti-Asian bias predicts interstate migration among second-generation Asian Americans with two Asian-born parents (Table~\ref{regtab-mig-01}). I find no significant correlation between bias and migration decisions. Among movers, self-identified Asians live in states with 0.02 standard deviations more bias than their birth states (standard error 0.04), a difference not distinguishable from zero; the positive sign means that selective migration, if anything, works against finding an effect of bias on identity switching.

In the CPS, a household respondent reports racial identity for all members, raising the concern that proxy responses may not reflect children's own identities. Two pieces of evidence mitigate this concern. First, \textcite{duncanIntermarriageIntergenerationalTransmission2011} show that reported Hispanic identity does not vary with respondent identity, and Table~\ref{tab:hispbyproxy} shows a similar pattern for Asian identity.\footnote{According to CPS guidelines, household respondents must be at least 15 years old with sufficient household knowledge. When the proxy is ``self,'' respondents range from 15 to 17 years old.} Second, I find similar patterns of ethnic attrition and bias effects among adult Asian Americans who self-report their own identity (Table~\ref{tab:asianbygen-adults}).

Finally, the main specification includes region $\times$ year fixed effects and identifies the effect of bias from differences across states within a region and year and from changes within states over time. It does not include state fixed effects, and this is a limitation. Appendix Table~\ref{regtab-all-fe} re-estimates the pooled linear probability model with region, state, year, and region $\times$ year fixed effects in every combination. Adding state fixed effects to the baseline specification (column 8) moves the bias coefficient from $-$0.10 (standard error 0.03) to $-$0.04 with a standard error of 0.03, and the specification with state and year fixed effects (column 4) gives $-$0.02 with the same standard error. Neither estimate is statistically distinguishable from zero. The 95 percent confidence interval in column 8, roughly $-$0.10 to 0.02, reaches the baseline estimate, while the interval in column 4, roughly $-$0.08 to 0.04, excludes it; the point estimates are less than half and about a fifth of the baseline. The imprecision does not come from a lack of within-state variation (the residual standard deviation of the index after state and region $\times$ year effects is 0.18, against 0.22 after region $\times$ year effects alone, and each state contributes at most five distinct values of the index); it is the point estimate itself that falls once state effects are included. Two explanations are consistent with that pattern and the data do not distinguish them: a fixed characteristic of states with high anti-Asian bias, rather than the bias itself, could lower Asian identification among children of mixed ancestry; or measurement error in the proxy-based index, which is proportionally larger for changes within a state than for differences between states, could attenuate the within-state estimate. A specification with state fixed effects and no time controls (column 3) yields a precise $-$0.02, but it attributes to bias any common time trend in identification and is not the baseline model with state effects added. The within-state estimates are therefore uninformative rather than reassuring, and the results in this paper should be read with the possibility of confounding by fixed state characteristics left open.

Larger Asian American populations could themselves affect state-level bias. Figure~\ref{scatter-plot-1} shows no correlation between state Asian American population shares and average anti-Asian bias.

\section{Conclusion}\label{sec:conc}

As American society becomes increasingly multiracial, racial identity choices will shape political representation, resource allocation, and social cohesion, and the statistics that policy relies on will continue to count people by the race they report. Understanding the determinants of identity is therefore important both for researchers studying discrimination's role in racial economic gaps and for the agencies that enforce civil rights law and monitor disparities with self-reported data. This paper shows how individual characteristics and anti-Asian sentiment influence racial identity reporting among Asian Americans.
 
Using multinomial logit analysis, I show that anti-Asian bias reshapes racial identity choices, with effects that vary across family structures and generational distance from immigration. Among pooled second-generation children of mixed-race parents, a one standard deviation increase in bias raises ``White only'' reporting by 12 percentage points and reduces ``Asian only'' by 10~pp. The effects differ by parental composition: in AW families, the same increase in bias reduces ``Asian only'' by 13~pp and increases ``White only'' by 20~pp, while in WA families, ``Asian only'' decreases by 12~pp and ``White only'' increases by 10~pp. Both between-model tests and a pooled model that interacts bias with an AW indicator confirm that the bias effect on ``White only'' identity is larger in AW families (by 10.6~pp in the between-model comparison and 15.1~pp in the pooled model). Among third-generation respondents, bias effects scale with the number of Asian grandparents: those with three Asian grandparents show the largest shifts ($+$19~pp ``White only''), while those with one grandparent show small, imprecise effects.
 
Parental education shapes identity choices in ways that also differ across family structures. In AW families, maternal college education increases ``Asian only'' by 10~pp and decreases ``White only'' by 12~pp, while in WA families maternal education produces no significant change---a difference the between-model tests find significant on both outcomes, although those tests do not account for state-level dependence. Among the third generation, parental education effects increase with ancestry concentration: maternal college raises ``Asian only'' by 32~pp for those with three Asian grandparents versus 7~pp for those with two, differences that are statistically significant by the same bootstrap criterion. However, the relationship between maternal education and identification varies by environment. For Asian-White (AW) families, the interaction between state-level bias and maternal college education is $-$0.23 (standard error 0.10), indicating that the negative relationship between bias and Asian identification is more pronounced among individuals with higher maternal education.
 
These results have consequences for the measurement of racial disparities and for the policies that depend on it. Because the combination of higher education and higher bias is associated with less, not more, Asian identification, the Asian-identified population in high-bias states is negatively selected on education and socioeconomic status. Outcomes measured on self-reported race therefore understate the outcomes of people of Asian ancestry in the most prejudiced areas, and comparisons of Asian and White outcomes there are biased in a direction that depends on which group is ahead: where Asian Americans trail, the gap is overstated, and where they lead, the lead is understated. Federal race and ethnicity statistics, hate crime reporting, and the disparity measures used to target and evaluate programs all rest on self-identification \autocite{ombRevisionsStatisticalPolicy2024, covid19HateCrimesAct2021}, so these distortions are not confined to academic estimates. More broadly, the strategic nature of identity decisions means that studies relying on self-reported racial identity may understate Asian American socioeconomic success in high-bias areas and understate the group's progress across generations.
 
One limitation qualifies these conclusions. The estimates rely on variation in bias across states within a region and year and on changes within states over time; when state fixed effects are added to the baseline specification, the estimates become too imprecise to be informative. The possibility that fixed characteristics of high-bias states, rather than bias itself, lower Asian identification among children of mixed ancestry therefore cannot be ruled out with these data.
 
This study encourages further research into relationships between bias and self-reported identities for other groups, including other ethnic and racial minorities and sexual minorities. Scholars could examine how recent anti-Asian violence following the COVID-19 pandemic has influenced identity patterns, providing natural experiments in bias effects. Researchers might also explore identity choices in specific institutional contexts like college admissions or workplace advancement, where model minority stereotypes create complex incentive structures, or extend the analysis to subgroups facing distinct stereotypes and discrimination patterns. More granular geographic analysis at the county or zip-code level could provide more precise measures of local bias, though data limitations prevent this in the current analysis. Understanding strategic racial identification among Asian Americans is essential for designing effective anti-discrimination policies and accurately measuring racial equity progress.
